\chapter{Аналитический раздел}

\section{Анализ предметной области}

Предметной областью данного проекта является видеохостинг-платформа, в которой сочетаются две основные стороны работы системы: просмотр контента и его публикация.
С одной стороны, пользователь пользуется платформой как зритель: ищет и просматривает интересующие видео.
С другой стороны, пользователь может выступать как автор, то есть загружать собственные видео в общий доступ.
Таким образом, в базе данных требуется хранить не только сведения о самих видеозаписях, но и связи между пользователями, каналами и действиями, которые выполняются в системе.

С точки зрения структуры данных, основные компоненты предметной области включают:
\begin{itemize}
	\item пользователей, которые могут просматривать или публиковать видео;
	\item каналы, создаваемые пользователями и содержащие видеозаписи;
	\item видео, доступные для взаимодействия.
\end{itemize}

\section{Анализ существующих решений}

Существуют готовые решения, предоставляющие пользователям возможность просмотра и публикации видеоконтента. Для определения требований к приложению рассмотрены три платформы:
\begin{itemize}
	\item VK Video;
	\item RuTube;
	\item Яндекс Дзен.
\end{itemize}
Анализ проведён в сравнительном формате по функциональным критериям, которые напрямую влияют на структуру приложения. Результаты анализа представлены в таблице~\ref{tab:solutions}.

\begin{longtable}{|p{7cm}|p{2.5cm}|c|c|}
\caption{Сравнение существующих решений предметной области}\label{tab:solutions} \\
\hline
\textbf{Критерий} & \textbf{VK Video} & \textbf{RuTube} & \textbf{Яндекс Дзен} \\
\hline
\endfirsthead
\multicolumn{4}{l}{\textit{Продолжение таблицы~\ref{tab:solutions}}} \\
\hline
\textbf{Критерий} & \textbf{VK Video} & \textbf{RuTube} & \textbf{Яндекс Дзен} \\
\hline
\endhead
\hline
\endfoot
\hline
\endlastfoot
Подписка на каналы & Да & Да & Да \\
\hline
Уведомления о новых видео & Да & Да & Да \\
\hline
Комментирование видео & Да & Да & Да \\
\hline
Формирование плейлистов & Да & Да & Нет \\
\hline
Оценивание видео & Частично & Частично & Частично \\
\hline
Оценивание комментариев & Да & Частично & Да \\
\hline
Владение несколькими каналами & Нет & Нет & Нет \\
\hline
\end{longtable}


\section*{Выводы по разделу}

%Проведенный анализ показывает, что базовый набор функций видеохостинга включает учет пользователей, хранение каналов,
%видеозаписей, подписок, комментариев и истории просмотров. Для проектируемой системы обязательными являются механизмы
%публикации видео и взаимодействия пользователей с контентом. Дополнительно целесообразно предусмотреть поддержку плейлистов,
%так как эта функция представлена в специализированных видеоплатформах и важна для структурирования контента.


\clearpage
